\section{System Architecture}\label{sec:arch}

The Vecna framework is constructed as a modular, decoupled architecture partitioned into four principal layers (Figure~\ref{fig:architecture}): data ingestion, heterogeneous feature extraction, high-throughput hybrid indexing, and multi-modal retrieval. By decoupling the offline processing pipeline from the online interactive search subsystem, the architecture ensures deterministic scalability while maintaining low-latency responses for interactive queries.

\begin{figure}[t]
\centering
\begin{tikzpicture}[
    node distance=0.28cm and 0.4cm,
    box/.style={rectangle, draw=blue!70!black, fill=blue!5, rounded corners=2pt, minimum width=3.3cm, minimum height=0.5cm, align=center, font=\scriptsize\sffamily},
    modelbox/.style={rectangle, draw=teal!70!black, fill=teal!8, rounded corners=2pt, minimum width=3.3cm, minimum height=0.5cm, align=center, font=\scriptsize\sffamily},
    dbbox/.style={cylinder, shape border rotate=90, draw=purple!70!black, fill=purple!5, aspect=0.18, minimum width=3.3cm, minimum height=0.65cm, align=center, font=\scriptsize\sffamily},
    procbox/.style={rectangle, draw=orange!80!black, fill=orange!8, rounded corners=2pt, minimum width=3.3cm, minimum height=0.5cm, align=center, font=\scriptsize\sffamily},
    arrow/.style={->, >=stealth, thick, draw=gray!70!black}
]

% Offline Column
\node[box, fill=gray!10, draw=gray!60] (video) {Raw Video Archive\\(I-Frames + $\Delta t_{\max} = 2.0\text{s}$)};
\node[modelbox, below=0.26cm of video] (models) {Multi-Model Feature Extraction\\OpenCLIP, SigLIP, Qwen3-VL, WhisperX, OCR};
\node[dbbox, below=0.26cm of models] (milvus) {Milvus Vector Database\\SCANN (Dense) \& BM25 (Sparse)};

% Online Query Flow
\node[box, fill=yellow!10, draw=yellow!60!black, below=0.30cm of milvus] (query) {User Query (Text, Tags, Multi-Step Splits)};
\node[procbox, below=0.26cm of query] (hyde) {Factual HyDE Expansion (Groq LLM)};
\node[procbox, below=0.26cm of hyde] (search) {Multi-Vector \& BM25 Hybrid Retrieval};
\node[procbox, below=0.26cm of search] (cluster) {Segment Clustering \& Keyframe Diversification};
\node[procbox, below=0.26cm of cluster] (temporal) {DP Temporal Sequence Search ($\Delta t_{\max}$)};
\node[box, fill=green!10, draw=green!60!black, below=0.26cm of temporal] (ui) {Interactive Web UI (Timeline Nav \& Target Selection)};

% Arrows
\draw[arrow] (video) -- (models);
\draw[arrow] (models) -- (milvus);
\draw[arrow] (query) -- (hyde);
\draw[arrow] (hyde) -- (search);
\draw[arrow] (milvus.east) -- ++(0.35,0) |- (search.east);
\draw[arrow] (search) -- (cluster);
\draw[arrow] (cluster) -- (temporal);
\draw[arrow] (temporal) -- (ui);

\end{tikzpicture}
\vspace{-2mm}
\Description{System architecture and operational workflow of the Vecna multimodal video retrieval framework.}
\caption{System architecture and operational workflow, unifying multi-model extraction, Milvus SCANN/BM25 indexing, segment diversification, and temporal sequence search.}
\label{fig:architecture}
\vspace{-2mm}
\end{figure}

\subsection{Adaptive I-Frame Keyframe Selection}
To efficiently process extensive video corpora without redundant computational overhead, Vecna eschews naive fixed-interval sampling in favor of an adaptive keyframe selection mechanism predicated on intrinsic video encoding properties. Utilizing \texttt{ffprobe} packet flags, we isolate native Intra-coded frames (I-frames), which demarcate definitive scene cuts or group-of-pictures (GOP) boundaries. To guarantee continuous temporal coverage across static or low-motion scenes with sparse I-frames, the ingestion pipeline enforces a maximum temporal gap ceiling $\Delta t_{\max} = 2.0$ seconds ($\text{max\_scene\_length} = 2.0\text{s} \times \text{FPS}$):
\begin{equation}
    t(k_{i+1}) - t(k_i) \leq \Delta t_{\max}
\end{equation}
If the temporal gap between consecutive keyframes reaches $\Delta t_{\max}$, a new keyframe is automatically sampled. Crucially, keyframes are stored at native $1280 \times 720$ resolution for fine-grained verification, while compressed thumbnails at $0.25\times$ scale ($320 \times 180$) are pre-generated to facilitate sub-second grid rendering in the interactive interface.

\subsection{Heterogeneous Feature Extraction}
The representational capacity of the framework hinges on a suite of complementary extraction backends designed to encapsulate distinct modalities:
\begin{itemize}
    \item \textbf{Visual Semantics:} Keyframes are simultaneously mapped into dense embedding spaces using three foundational visual encoders: OpenCLIP \texttt{PE-Core-L-14-336} ($1024$-dimensional normalized visual vectors from $336 \times 336$ crops using the meta pretrained checkpoint, excelling at fine-grained textures and small objects) \cite{radford2021learning, ilharco2021openclip}, Google SigLIP \texttt{ViT-SO400M-14-SigLIP-384} ($1152$-dimensional embeddings trained with pairwise sigmoid loss on WebLI, providing robust zero-shot generalization for human actions and scene context) \cite{zhai2023sigmoid}, and Qwen3-VL Embedding \texttt{Qwen3-VL-2B} ($2048$-dimensional multimodal vectors capturing complex visual compositions and semantic relations) \cite{bai2023qwen}.
    \item \textbf{Optical Character Recognition (OCR):} Textual elements localized within video frames are extracted through a two-stage pipeline: text detection via PaddleOCR (optimized with a dedicated Vietnamese lexical dictionary) \cite{du2020ppocr} and sequence recognition via VietOCR (\texttt{vgg\_transformer}), robustly extracting on-screen text overlays, news tickers, and signboards into searchable text strings. Extracted text is subsequently embedded into $1024$-dimensional vectors utilizing BGE-M3 \cite{chen2024bge}.
    \item \textbf{Automatic Speech Recognition (ASR):} Audio streams are demuxed and segmented via Voice Activity Detection (VAD) utilizing \texttt{PyAnnote.audio}, followed by high-speed transcription and phoneme-level word alignment using WhisperX Turbo (\texttt{large-v3-turbo}) \cite{bainix2023whisperx}. Transcribed linguistic tokens are aligned with keyframe timestamps and embedded via BGE-M3.
    \item \textbf{Canonical Frame Identification:} Extracted visual, OCR, and speech representations are uniquely bound to canonical frame identifiers formatted as \texttt{VideoID\#FrameIndex} and ingested into Milvus collections.
\end{itemize}

\subsection{Overlapped ONNX DirectML Pipeline}
To maximize hardware utilization on heterogeneous Windows environments, the BGE-M3 dense embedding generation employs a highly optimized, three-stage asynchronous pipeline facilitated by ONNX Runtime and the DirectML execution provider. The workload is partitioned into concurrent threads: a CPU-bound Tokenizer thread, a GPU-bound DirectML Inference thread, and a persistent Disk Writer thread.

To mitigate out-of-memory (OOM) faults while maximizing throughput, the inference stage employs a dynamic batch sizing algorithm that adapts to varying token sequence lengths. The instantaneous batch size \( B(L) \) for a given sequence length \( L \) is computed as:
\begin{equation}
    B(L) = \min\left(B_{\max}, \left\lfloor B_{\text{base}} \cdot \left(\frac{L_{\text{ref}}}{L}\right)^2 \right\rfloor\right)
\end{equation}
where the absolute maximum batch size \( B_{\max} = 128 \), and the reference length \( L_{\text{ref}} = 256 \). This quadratic decay ensures optimal occupancy of the VRAM buffers independent of the input distribution.

\subsection{Hybrid Milvus Indexing}
The terminal stage of the offline pipeline constitutes the orchestration of the heterogeneous embeddings into a distributed vector database, specifically Milvus \cite{wang2021milvus}. To facilitate expeditious nearest-neighbor queries, we deploy a hybrid indexing schema:
\begin{itemize}
    \item \textbf{Dense Visual Embeddings:} High-dimensional visual features are partitioned utilizing the SCANN (Scalable Nearest Neighbors) index \cite{guo2020accelerating}, employing cosine similarity. The clustering topology is optimized with \( \text{nlist} = 512 \) partitions and a conservative \( \text{nprobe} = 8 \) parameter during traversal.
    \item \textbf{Sparse Textual Features:} The textual semantics derived from OCR and ASR pipelines are organized via a Term-at-a-Time (TAAT) inverted index leveraging the BM25 probabilistic retrieval framework, with parameters meticulously calibrated to \( k_1 = 1.2 \) and \( b = 0.75 \).
\end{itemize}
This hybrid architecture intrinsically supports the fusion of dense continuous vector representations and sparse discrete token occurrences, a critical requirement for open-domain interactive retrieval.
